\section{So what? Why any of this should matter to someone who isn't grading it}

Everything up to this point has been building a model and testing it. This section is me stepping back and asking the question I'd actually be asked by a risk manager, an investor, or an interviewer: fine, you built a working distance-to-default pipeline and it predicted Yes Bank's stress eighteen months early — so what should anyone actually \emph{do} differently because of this? I think there are real, specific answers, for several different readers, and I want to give each of them their own answer rather than one vague conclusion.

\subsection{For a bank investor or a bond holder}
If I had been holding Yes Bank's Tier~II bonds or shares in mid-2018, my model's own numbers say I had a genuinely usable early-warning signal available to me, for free, from public price data, starting in September 2018 — more than a year before the first rating agency action that seriously questioned Yes Bank's capital position. The honest caveat, which I don't want to undersell, is in my own robustness check: a DD-based alarm would also have fired falsely, in the sense that Yes Bank survived (as an independent-ish entity with new management) for over a year after the alarm first went off, and partially recovered in early 2019 without failing. A real investor using this signal in real time would have needed to decide, in September 2018, whether a four-standard-deviation collapse in one month was a reason to exit entirely or a reason to just reduce exposure and watch closely — and my model, by itself, doesn't tell you which. What it does tell you is that waiting for a rating agency to confirm the story before acting means acting, on average across my three measured episodes, somewhere between three days and over two months later than the market had already told you.

\subsection{For a bank's own risk management desk}
This is where I think the project's biggest practical lesson sits, and it isn't ``use Merton instead of ratings.'' It's that the two signals fail in \emph{opposite, complementary} ways, which is exactly why a real risk desk should want both running simultaneously, not one or the other. My DD series is loud, immediate, and chronically pessimistic once a firm is under real stress — exactly the kind of signal you want for day-to-day monitoring, precisely because it updates every single trading session and costs nothing to compute beyond a share price. But it is structurally blind to anything that happens outside the market's view: it could not see the December-2019 loss before 13 March 2020 because the market itself couldn't see it. A rating agency, for all its slowness, has analysts who can sit across a table from management, ask for internal numbers the market never sees, and — this is the part I think gets underrated in most critiques of agencies, including the one I opened this whole project with — occasionally catch something the market hasn't priced in yet \emph{because it isn't public}. A risk desk that only watched DD would have missed nothing about Yes Bank's ongoing distress but would have been just as surprised by the specific timing of the moratorium as everyone else. A risk desk that only watched ratings would have been months late to the party entirely. Running both, and specifically watching for the moments they \emph{disagree}, is the actual operational lesson I'd pull from this project if I were building a monitoring dashboard rather than writing a note about one.

\subsection{For a regulator, and for Indian financial regulation specifically}
I think my results are an argument, even if a small, single-case one, for why the RBI increasingly leans on its own direct, non-public supervisory channels — on-site inspections, divergence audits, confidential reporting — rather than relying on either market prices or published ratings as its primary early-warning tool. Both of the public signals I tested were, in their own ways, informative but insufficient: the market was fast but blind to what it couldn't see, and the agencies were seeing much of the same public information the market saw, just slower. The one genuinely decisive, non-public piece of information in this entire case — Yes Bank's real December-2019 numbers — only became visible to the RBI through its own supervisory access, not through anything a market-based model or a public rating could have picked up in time. If I were advising a regulator rather than a student writing a project, my takeaway wouldn't be ``trust market signals more.'' It would be ``market signals are a genuinely useful early, cheap, continuously-updating tripwire for \emph{where} to point your scarce supervisory attention — they are not a substitute for having supervisory access in the first place.''

\subsection{The specific, structural problem with the rating agencies, now that I've actually checked it}
When I started this project, my working assumption — shaped partly by the general global critique in Section~\ref{sec:critique} and partly, I'll admit, by something closer to received wisdom than verified fact — was that the rating agencies had probably been pressured into going slow on Yes Bank specifically. Having now actually read SEBI's 2023 order, I think the real, verified picture is more useful than the simpler story I started with, and it's worth spelling out why. SEBI's forensic audit \emph{did} find that CARE's senior management interfered with and delayed a rating action — for DHFL, a different, related stressed borrower, in the same period. For Yes Bank specifically, the same investigation, after years of cross-examination, concluded the interference charge couldn't be sustained on the evidence available. I think the honest, structural lesson here is not ``CARE was corrupt about Yes Bank'' (the regulator looked hard and didn't find enough to prove that) but something a little more unsettling precisely because it's more mundane: \textbf{the same issuer-pays conflict of interest that was proven to have delayed one bank's downgrade was structurally present, and plausible, across the whole stressed-lender portfolio of that period} — DHFL, IL\&FS, and Yes Bank were all being rated by the same analysts, under the same commercial pressure, in the same eighteen-month window, and the one case SEBI could actually prove is not obviously a special case rather than a sample. A model like mine doesn't have this problem at all, for a very simple, almost dull reason: a share price is set by thousands of independent, anonymous traders with money at risk, not by a committee whose employer is paid by the entity being graded. That structural independence — not any claim about speed or sophistication — is, I now think, the real argument for keeping a market-based signal running alongside agency ratings, and it's a stronger and more specific argument than the one I'd have made before I did the verification work in Section~\ref{sec:case}.

\subsection{What I'd tell my own FRM-prep self, a year ago}
I'm building this project partly as an FRM candidate, so I want to end this section on the question that's most practical for me personally: what does this change about how I'd answer a credit-risk question on an exam, or in an interview? Three things, concretely. First, I now actually understand \emph{why} distance-to-default is quoted in standard-deviation units rather than rupees or percentages — because I had to solve the two-equation system myself and watch what happened to $\sigma_V$ as leverage changed, rather than just memorising the formula. Second, I now have a real, specific, defensible answer to ``what's wrong with structural models for banks'' that goes beyond the textbook line about deposits being hard to model — I have an actual month (October 2019) where I watched $d_2$ go negative in my own data and had to sit with what that number was and wasn't telling me. Third, and this is the one I didn't expect going in: I now think the single most exam-relevant, interview-relevant fact from this whole project is the March 2020 whiplash itself — a model can be \emph{correctly} reporting that priced risk has fallen while the \emph{true} risk is rising, because the two things aren't the same, and the gap between them is exactly where supervisory information earns its keep. That's a sentence I could only have written after watching my own pipeline produce a number that surprised me.
